\section{System Architecture}

\subsection{Node Topology}

The framework operates as a distributed processing architecture
partitioned across two functionally isolated computational nodes.
The primary \textit{Gait Generator} node manages all rhythmic
trajectory computation, inverse kinematics resolution, and
actuator command dispatch. An asynchronous \textit{Balance
Controller} node exclusively handles inertial measurement
acquisition, signal conditioning, and PID correction computation.
Communication between the two nodes is strictly limited to
lightweight corrective vector messages. This architectural
segregation guarantees that trajectory generation anomalies
cannot preempt or block critical sensory compensations, and
vice versa.

\subsection{State Representation}

Global governance overlying the control architecture is dictated
by boolean gating, handling two discrete operating states:

\begin{itemize}
    \item \textbf{HOME:} A static configuration lacking cyclic
    translation vectors. The controller computes uniform
    adjustments in joint space rigidly mapped to a resting stance.
    Base angular deviations translate proportionally into
    antagonistic torque compensations aimed at continuous postural
    rectification. A configurable startup delay of $T_{delay}$
    seconds prevents premature correction during system
    initialization transients.

    \item \textbf{GAIT:} The dynamic state active during commanded
    planar translations. The platform executes cyclic polynomial
    trajectories using configurable phase shift arrays to generate
    trot ($\phi = [0, 0, 0.5, 0.5]$), walk ($\phi = [0, 0.25,
    0.5, 0.75]$), and other gaits involving multiple limbs.
    Spatial corrections derived from inertial feedback are applied
    exclusively to limbs in stance phase, while swing legs maintain
    zero compensation to ensure uninhibited stride clearance.
\end{itemize}
